\documentclass[final,t]{beamer}

\usepackage[size=a0,orientation=portrait,scale=1.00]{beamerposter}
\usepackage[utf8]{inputenc}
\usepackage[T1]{fontenc}
\usepackage{lmodern}
\usepackage[spanish,es-tabla]{babel}
\usepackage[table]{xcolor}
\usepackage{graphicx}
\usepackage{amsmath}
\usepackage{booktabs}
\usepackage{tabularx}
\usepackage{array}
\usepackage{enumitem}
\usepackage{tikz}
\usetikzlibrary{positioning,calc}
\graphicspath{{./}{../}}

\setbeamertemplate{navigation symbols}{}
\setbeamertemplate{headline}{}
\setbeamertemplate{footline}{}
\setbeamersize{text margin left=14mm,text margin right=14mm}
\renewcommand{\familydefault}{\sfdefault}
\setlength{\parindent}{0pt}
\setlength{\parskip}{0.20em}
\renewcommand{\arraystretch}{1.12}
\newcolumntype{Y}{>{\raggedright\arraybackslash}X}

\definecolor{paperText}{HTML}{1A1A1A}
\definecolor{paperMuted}{HTML}{5F6368}
\definecolor{paperRule}{HTML}{B7B7B7}
\definecolor{paperAccent}{HTML}{1F4E79}
\definecolor{paperAccentLight}{HTML}{EAF2F8}
\definecolor{barOne}{HTML}{315C8C}
\definecolor{barTwo}{HTML}{5F84B2}
\definecolor{barThree}{HTML}{8EABCC}
\definecolor{barFour}{HTML}{B5C7DC}
\definecolor{barOther}{HTML}{D9E2EC}

\setbeamercolor{background canvas}{bg=white}
\setbeamercolor{normal text}{fg=paperText,bg=white}
\setbeamerfont{normal text}{size=\large}

\newcommand{\articleblock}[2]{%
  {\color{paperAccent}\LARGE\bfseries #1}\par
  \vspace{0.05cm}{\color{paperRule}\hrule height 1.0pt}\vspace{0.08cm}
  {#2}\par\vspace{0.08cm}
}
\newcommand{\smallnote}[1]{{\color{paperMuted}\normalsize #1}}

\newcommand{\metriccell}[3]{%
  \begin{minipage}[t]{0.31\linewidth}
    \centering
    {\fontsize{58}{62}\selectfont\bfseries #1}\par
    {\Large\bfseries #2}\par
    \smallnote{#3}
  \end{minipage}%
}

\newcommand{\funnelrow}[5]{%
  \noindent\begin{minipage}[t]{0.30\linewidth}\raggedright\Large #1\end{minipage}%
  \hfill
  \begin{minipage}[t]{0.56\linewidth}
    \vspace{0.20em}%
    \begin{tikzpicture}[x=1cm,y=1cm]
      \pgfmathsetmacro{\barwidth}{#3*(#2)/100.0}%
      \draw[fill=#5, draw=none] (0,0) rectangle (\barwidth,0.50);
      \draw[draw=paperRule, line width=0.35pt] (0,0) rectangle (#3,0.50);
    \end{tikzpicture}%
  \end{minipage}%
  \hfill
  \begin{minipage}[t]{0.11\linewidth}\raggedleft\Large\bfseries #4\end{minipage}\par\vspace{0.14cm}%
}

\begin{document}
\begin{frame}[t]

\vspace*{0.12cm}
\begin{center}
  
\includegraphics[height=3.0cm]{ufro-logo.png}\par
  \vspace{0.15cm}
  {\fontsize{50}{56}\selectfont\bfseries Taxonomía de motivos de retrabajo en solicitudes de cambios de agentes de IA}\par
  \vspace{0.15cm}
  {\Large Javier Alcalde y Diego Labrín \quad | \quad Profesor: Dr. Pablo Valenzuela}\par
  \vspace{0.05cm}
  {\large Departamento de Computación e Informática, Universidad de La Frontera \quad | \quad Minería de Repositorios de Software}\par
  \vspace{0.12cm}
  {\color{paperRule}\rule{0.94\textwidth}{1.5pt}}
\end{center}

\vspace{0.15cm}
\begin{center}
\begin{minipage}{0.94\textwidth}

\articleblock{Pregunta y objetivos}{
Los agentes de IA pueden producir solicitudes integrables, pero la integración requiere retrabajo humano. Analizamos solicitudes del dataset \textbf{AIDev} (33.596 solicitudes generadas por IA en repositorios reales) \textbf{integradas después de retrabajo observable}.

\vspace{0.12cm}
\begin{tabularx}{\linewidth}{>{\bfseries}p{5.5cm}Y}
Pregunta & ¿Qué motivos de retrabajo humano emergen antes de la integración? \\
Obj. 1 & Definir población operacional auditable desde AIDev. \\
Obj. 2 & Implementar muestra estratificada reproducible por agente. \\
Obj. 3 & Generar tarjetas con evidencia y aplicar clasificación abierta. \\
\end{tabularx}
}

\articleblock{Metodología}{
\begin{center}
\begin{tikzpicture}[x=1cm,y=1cm]
  \node[draw=paperAccent, fill=paperAccentLight, rounded corners=5pt, minimum width=14.0cm, minimum height=1.8cm, align=center] (prep) at (0,0) {\bfseries Preparación\\\smallnote{filtros, muestra, tarjetas}};
  \node[draw=paperAccent, fill=white, rounded corners=5pt, minimum width=14.0cm, minimum height=1.8cm, align=center] (exec) at (16.5,0) {\bfseries Ejecución\\\smallnote{clasificación abierta}};
  \node[draw=paperAccent, fill=paperAccentLight, rounded corners=5pt, minimum width=14.0cm, minimum height=1.8cm, align=center] (ana) at (33.0,0) {\bfseries Análisis\\\smallnote{taxonomía, métricas}};
  \draw[-latex, line width=1.3pt, paperAccent] (prep) -- (exec);
  \draw[-latex, line width=1.3pt, paperAccent] (exec) -- (ana);
\end{tikzpicture}
\end{center}

\vspace{0.10cm}
{\small\textbf{Confiabilidad:} 2 evaluadores, 300 tarjetas. Acuerdo categoría padre: 73\%; subcategoría: 47\%.}

\vspace{0.12cm}
\begin{tabularx}{\linewidth}{>{\bfseries}p{4.2cm}Y>{\bfseries}p{4.2cm}Y}
Inclusión & Integrada, \texttt{commit\_count>1}, \texttt{human\_comment\_count>0}. & Evidencia & Comentarios, revisiones, timeline. \\
Exclusión & Abiertas, sin integrar, sin comentarios humanos. & Validez & Cada asignación responde la pregunta. \\
\end{tabularx}
}

\articleblock{Datos y muestra}{
\begin{center}
\metriccell{33.596}{Universo}{AIDev}\hfill
\metriccell{3.166}{Población}{operacional}\hfill
\metriccell{300}{Muestra}{estratificada}
\end{center}

\vspace{0.25cm}
\funnelrow{Universo AIDev}{100}{31.0}{100\%}{barOne}
\funnelrow{Integradas}{71.48}{31.0}{71\%}{barTwo}
\funnelrow{Con commits}{20.49}{31.0}{20\%}{barThree}
\funnelrow{Población}{9.42}{31.0}{9\%}{barFour}
\funnelrow{Muestra}{0.89}{31.0}{1\%}{barOther}

\vspace{0.04cm}
{\small\color{paperMuted}Error muestral: $\pm 5{,}4\%$ (95\% IC). Cuotas: Copilot 145, Devin 86, Codex 45, Cursor 17, Claude 7.}
}

\articleblock{Resultado: taxonomía final}{
Las 4 primeras familias: \textbf{228 de 300} (76\%). El treemap muestra frecuencia de motivos de retrabajo pre-integración.

\vspace{0.08cm}
\begin{center}
\includegraphics[width=0.90\linewidth]{taxonomia-final-merged-after-rework-treemap.pdf}
\end{center}
}

\articleblock{Discusión}{
\small
\textbf{P1 (34,7\%):} Confirma Liang et al.~[1] (38\%). Mejora en validación, pero persiste como barrera principal.

\textbf{P2 (19,0\%):} Vaithilingam et al.~[2] (22\% en novatos). IA introduce errores sutiles incluso con revisión.

\textbf{P3 (12,3\%):} Barke et al.~[3] ``sobrecódigo generado''. Cuantificado en contexto real.

\textbf{P4 (10,0\%):} Hallazgo emergente: IA debe navegar reglas sociales del repositorio.
}

\articleblock{Implicancias}{
\small
\begin{itemize}[leftmargin=*, label=\textcolor{paperAccent}{\textbullet}, itemsep=0.06cm]
\item \textbf{Desarrolladores:} Ejecutar verificaciones locales (CI, tests, lint) antes de solicitar revisión; usar taxonomía como lista de verificación pre-envío.
\item \textbf{Investigadores:} Medir trayectoria de integración (confirmaciones, cambios, tiempo), no solo tasa de integración; autonomía real es menor que la reportada.
\item \textbf{Constructores de Herramientas:} Incorporar memoria de repositorio (convenciones, responsables, dependencias); bucles de auto-verificación antes de pedir revisión.
\end{itemize}
}

\articleblock{Alcance y límites}{
\small
\begin{tabularx}{\linewidth}{>{\bfseries}p{3.5cm}YY}
\toprule
Dimensión & Estado & Precaución \\
\midrule
Hecho & Población (3.166), muestra (n=300), taxonomía trazable. & Motivo principal; puede haber múltiples señales. \\
RQ2 (futuro) & Tiempo hasta integración por categoría. & Sin conclusiones temporales aún. \\
\bottomrule
\end{tabularx}

\vspace{0.08cm}
{\footnotesize\textbf{Reproducibilidad:} Scripts versionados, semilla \texttt{20260510}. \textbf{Repo:} \texttt{github.com/Diego-LC/proyecto-semestral-msr-aidev}.}
}

\articleblock{Referencias}{
\footnotesize
[1] P. Liang, Y. Yang, D. Myers, ``An Empirical Study of Quick Remedy Commits,'' \textit{Proc. ICSE}, 2024. doi: 10.1145/3597503.3623332

[2] P. Vaithilingam, T. Zhang, E. Glassman, ``Expectation vs. Experience: Evaluating the Usability of Code Generation Tools,'' \textit{Proc. CHI EA}, 2022. doi: 10.1145/3491102.3517585

[3] M. Barke, M. B. James, N. Polikarpova, ``Grounded Copilot: How Programmers Interact with Code-Generating Models,'' \textit{Proc. ACM Program. Lang.}, vol. 7, no. OOPSLA1, pp. 935--963, 2023. doi: 10.1145/3591241

[4] H. Li, ``AIDev: Dataset of AI-Generated Pull Requests,'' \textit{Hugging Face}, 2024. \texttt{https://huggingface.co/datasets/hao-li/AIDev}
}

\end{minipage}
\end{center}

\end{frame}
\end{document}